\documentclass{article}

\usepackage{mathtools}
\usepackage{amssymb}
\usepackage{amsthm}
\usepackage[spaceRequire = false]{algpseudocodex}
\usepackage{siunitx}
\usepackage{zref-clever}
\usepackage{../latex-course}
\usepackage{week-05}
\newcommand{\iter}[2]{#1_{(#2)}}           % the algorithms' iterates

\sisetup{per-mode = power}

\courseweek{05}{Theorems, Algorithms, Code, Units}
\coursestudents{Sandrine Pudukkattukaran Joy \and Lennard Boehm}

\title{The Gradient Descent Method}
\author{Sandrine Pudukkattukaran Joy (4743017), Lennard Boehm (4761574)}
\begin{document}
\maketitle

\section{Descent Methods}
\label{section:descent}
We study the minimization of an objective function $f \colon \R^n \to \R$, of which we assume the following.

\begin{assumption}
	\label{assumption:standing}
	The objective $f \colon \R^n \to \R$ is continuously differentiable, and $M \in \R^{n \times n}$ is symmetric and positive definite.
\end{assumption}

\begin{definition}[directional derivative, descent direction]
	\label{definition:descent}
	For $x, d \in \R^n$, the \emph{directional derivative} of $f$ at $x$ in direction $d$ is
	\begin{equation*}
		\dderiv{x}{d} \coloneqq \lim_{t \to 0} \frac{f(x + t \, d) - f(x)}{t}.
	\end{equation*}
	The direction $d$ is a \emph{descent direction} for $f$ at $x$ if $\dderiv{x}{d} < 0$.
\end{definition}

\begin{remark}
	Because $f$ is differentiable, the map $d \mapsto \dderiv{x}{d}$ is linear.
	This linear functional matches the \emph{derivative} $f'(x) \in (\R^n)^*$, so that we have $\dderiv{x}{d} = f'(x) \, d$.
\end{remark}
We call $x$ \emph{stationary} when $f'(x) = 0$, and where $f'(x) \neq 0$, the descent directions at $x$ form the open half-space $\{\, d \in \R^n : \dderiv{x}{d} < 0 \,\}$, bounded by the hyperplane on which $f'(x)$ vanishes.
A \emph{line-search descent method} starts from an initial guess $\iter{x}{0}$ and repeatedly picks a descent direction, picks a step length that produces descent along it, and moves, until a stopping criterion is met.
This core behavior is described in the following pseudocode.

\begin{algorithm} \hfill
\label{algorithm:descent}
	\begin{algorithmic}[1]
		\Procedure{DescentMethod}{$f, \iter{x}{0}$}
			\State $k \gets 0$
			\While{stopping criterion not satisfied}
		
			\State  choose descent direction $\iter{d}{k}$ 
				\label{algorithm:descent:direction}for $f$ at $\iter{x}{k}$
				
				\State  choose step length $\iter{t}{k} > 0$ with $f(\iter{x}{k} + \iter{t}{k} \, \iter{d}{k}) \le f(\iter{x}{k})$
				\label{algorithm:descent:steplength}
				\State $\iter{x}{k+1} \gets \iter{x}{k} + \iter{t}{k} \, \iter{d}{k}$
				\State $k \gets k + 1$
			\EndWhile
			\State \Return $\iter{x}{k}$
		\EndProcedure
	\end{algorithmic}
\end{algorithm}

Two important choices in \zcref{algorithm:descent} remain open: the search direction $\iter{d}{k}$ (\zcref{algorithm:descent:direction}) and the step length $\iter{t}{k}$ (\zcref{algorithm:descent:steplength}), which we address in the following section.

\section{Steepest-Descent Directions}
\label{section:direction}
The derivative $f'(x)$ contains all descent information for $f$, but it is an object in the dual space $(\R^n)^*$, whereas a search direction must be a vector in $\R^n$.
Fixing an inner product on the search space, we can find a best representative in the primal space, which makes descent information usable in the primal context.
For the symmetric positive definite matrix $M$ of \zcref{assumption:standing}, we set
\begin{equation*}
	\inner{u}{v}_M \coloneqq u^\transp M \, v
	\qquad\text{and}\qquad
	\norm{u}_M \coloneqq \inner{u}{u}_M^{1/2}.
\end{equation*}
Since $d \mapsto \dderiv{x}{d}$ is linear, there is a unique vector $\gradM f(x) \in \R^n$ with
\begin{equation}
	\label{eq:m-gradient}
	f'(x)\, d = \dderiv{x}{d} = \inner{\gradM f(x)}{d}_M \qquad \text{for all } d \in \R^n,
\end{equation}
called the \emph{$M$-gradient} of $f$ at $x$, and it can be computed by solving a linear system of equations.
A \emph{direction of steepest descent} at a non-stationary $x$ minimizes the normalized directional derivative $\dderiv{x}{d} / \norm{d}_M$ over $d \neq 0$.
By \eqref{eq:m-gradient} the numerator equals $\inner{\gradM f(x)}{d}_M$, so the Cauchy--Schwarz inequality shows the minimum is attained exactly at the negative multiples of $\gradM f(x)$.
Its negative is therefore a descent direction at every non-stationary point, since
\begin{equation}
	\label{eq:m-descent}
	\dderiv{x}{-\gradM f(x)} = \inner{\gradM f(x)}{-\gradM f(x)}_M = -\norm{\gradM f(x)}_M^2 < 0.
\end{equation}

\begin{remark}
	The gradient depends on the inner product, and so do its physical dimensions.
	Suppose $x$ is a position measured in $\unit{\meter}$ and $f$ an energy measured in $\unit{\joule}$, as for the potential $f(x) = m g \, x$ of a weight $m g = \qty{9.81}{\newton}$.
	Then $\dderiv{x}{d} = f'(x)\, d$ is an energy as well, so $f'(x)$ carries $\unit{\joule\per\meter} = \unit{\newton}$.
	For $\gradM f(x)$ to come out as a genuine displacement in $\unit{\meter}$, so that the update $\iter{x}{k+1} = \iter{x}{k} - t \, \gradM f(\iter{x}{k})$ runs with a dimensionless step length $t$, the matrix $M$ must carry the dimensions of a Hessian, i.e., $\unit{\joule\per\meter\squared}$.
\end{remark}

\section{Armijo Backtracking}
\label{section:armijo}
Having found descent directions that we can leverage in the descent algorithm, the line-search method still requires choosing a step length to follow a descent direction along for.
There are stalling issues to navigate in order to make sure the final algorithm converges, which leads to the \emph{Armijo condition} for a descent direction $d$ at $x$ and a parameter $\sigma \in (0, 1)$, which states that a step $t > 0$ is admissible when
\begin{equation}
	\label{eq:armijo}
		\frac{f(x + td) - f(x)}{t\dderiv{x}{d}} \ge \sigma.
\end{equation}
It formalizes the expectation that the actual objective decrease is at least a fixed fraction of the descent the first order model promised.
In practice, one decreases the step length until the condition is met.

\begin{algorithm}[Armijo backtracking] \hfill
	\label{algorithm:armijo}
	\begin{algorithmic}[1]
		\Require line-search function $t\mapsto f(x + t\,d)$
		\Require Armijo parameter $\sigma \in (0, 1)$, backtracking parameter $\beta \in (0, 1)$, initial step length $s > 0$
		\Ensure a step length $t > 0$ with sufficient descent
		\State Set $t \coloneqq s$
		\While{$f(x + t\,d) - f(x) > t\sigma\dderiv{x}{d}$}
			\State Set $t \coloneqq \beta \, t$
		\EndWhile
		\State \Return $t$
	\end{algorithmic}
\end{algorithm}

This is in fact a finite algorithm, as the following theorem shows.

\begin{theorem}[well-definedness of Armijo backtracking]
	\label{theorem:armijo}
	Let $\sigma \in (0, 1)$. For every $x, d \in \R^n$ with $\dderiv{x}{d} < 0$ there exists a $T > 0$ such that the Armijo condition \eqref{eq:armijo} holds for all $t \in (0, T]$.
\end{theorem}

\begin{proof}
	Suppose not. Then there is a sequence $\iter{t}{k} \searrow 0$ with $f(x + \iter{t}{k} d) > f(x) + \sigma \, \iter{t}{k} \, \dderiv{x}{d}$, hence
	\begin{equation*}
		\frac{f(x + \iter{t}{k} d) - f(x)}{\iter{t}{k}} > \sigma \, \dderiv{x}{d}.
	\end{equation*}
	Passing to the limit $k \rightarrow \infty$ gives $\dderiv{x}{d} \ge \sigma \, \dderiv{x}{d}$, which contradicts $\dderiv{x}{d} < 0$ because $\sigma \in (0, 1)$.
\end{proof}

Since the trial steps $s\beta^\ell$ tend to $0$, one of them eventually lies in $(0, T]$, so \zcref{algorithm:armijo} terminates whenever the search direction is a descent direction.

\section{\texorpdfstring{$M$}{M}-Gradient Descent}
\label{section:gradient}
The \emph{$M$-gradient descent method} is \zcref{algorithm:descent} with the search direction fixed to the negative $M$-gradient, the step length supplied by the Armijo backtracking rule of \zcref{section:armijo} and a termination criterion.

\begin{algorithm}[$M$-gradient descent with Armijo line search] \hfill
	\label{algorithm:gradient}
	\begin{algorithmic}[1]
		\Require objective $f \colon \R^n \to \R$
		\Require initial guess $\iter{x}{0} \in \R^n$
		\Require symmetric positive definite $M \in \R^{n \times n}$
		\Require Armijo parameter $\sigma \in (0, 1)$, backtracking parameter $\beta \in (0, 1)$, initial step length $s > 0$
		\Ensure a point $x \in \R^n$ with $f(x) \le f(\iter{x}{0})$
		\State Set $x \coloneqq \iter{x}{0}$
		\While {no stopping criterion satisfied}
			\State Compute $d = -\gradM f(x)$ from \eqref{eq:m-gradient}
			\State Determine a corresponding $t > 0$ by \zcref{algorithm:armijo}
			\State  Set $x \coloneqq x + t \, d$
		\EndWhile
		\State \Return $x$
	\end{algorithmic}
\end{algorithm}

At every non-stationary iterate, the direction $d = -\gradM f(x)$ is a descent direction by \eqref{eq:m-descent}, so $\dderiv{x}{d} < 0$ and the Armijo backtracking terminates by \zcref{theorem:armijo}; \zcref{algorithm:gradient} is therefore well defined.
\end{document}
